\phantomsection
\chapter{\lang{Case Study Context and Concepts}{Contexto do Estudo de Caso e Conceitos}}
\phantomsection

\section{libp2p: A Modular Networking Stack for Peer-to-Peer Applications}

Historicamente, aplicações \textit{peer-to-peer} (como Napster ou BitTorrent) possuíam camadas de rede monolíticas, com protocolos de transporte e mecanismos de descoberta de pares
rigidamente acoplados à lógica de suas aplicações. Isso criava sistemas inflexíveis, incapazes de se comunicar entre nodos usando diferentes meios
de transporte ou ambientes (navegador, servidor, smartphone).

A motivação por trás da libp2p foi criar uma \textit{stack} de rede modular agnóstica. A internet é um ambiente fora de controle, repleto de \textit{firewalls}, roteadores NAT e dispositivos
operando em camadas e contextos distintos. A libp2p foi imaginada para integrar esses diferentes sistemas, independentemente de seus protocolos, redes, sistemas operacionais e
ambientes distintos \cite{libp2p}. Para isso, ela é composta por uma coleção de módulos que podem ser empilhados e usados em conjunto, conforme a necessidade da aplicação.

Diferente do modelo cliente-servidor tradicional, onde os nós são identificados por endereços de IP, na libp2p cada nó é identificado por um PeerID, derivado da sua chave criptográfica pública \cite{libp2ppeerids}.
O endereço de IP é tratado apenas como uma rota temporária (multiaddr) para chegar àquela identidade \cite{multiaddr}.

A biblioteca suporta diversos meios de transporte, como TCP, QUIC (sobre UDP), WebRTC (para \textit{browsers}) e WebSockets \cite{libp2pspecs}. Uma vez que
dois nós tenham estabelecido uma conexão segura, é permitido que múltiplos subprotocolos paralelos rodem sobre essa mesma conexão (multiplexação), sem interferir uns nos outros \cite{libp2pconnections}. A ferramenta também
resolve o clássico problema de encontrar nós, através de módulos intercambiáveis \cite{libp2pspecs}.

Já para o problema de NAT, a libp2p oferece três módulos: o AutoNAT, que detecta se o nó está atrás de um NAT \cite{libp2pautonat}; o Circuit Relay, que usa um outro nó como intermediário
da conexão \cite{libp2prelay}; e o DCUtR, que tenta transformar essa conexão intermediada em uma conexão direta, através de \textit{hole punching} \cite{libp2pdcutr}.

A forte adoção da libp2p validou o modelo modular para toda uma indústria. Ela atua como infraestrutura de transporte em grandes redes como a camada de consenso do Ethereum \cite{ethp2p} e o IPFS \cite{libp2p}.

\section{Iroh}

Embora a libp2p e o IPFS (para sistemas de arquivos) tenham se estabelecido como padrões na indústria para redes \textit{peer-to-peer} e Web3, essas ferramentas ainda são consideradas, para algumas
aplicações, complexas. Como alternativa, surgiu uma segunda opção, que evoluiu conceitualmente e respondeu a gargalos da libp2p, com foco em velocidade, simplicidade e dispositivos móveis.
O acoplamento entre os componentes da libp2p (DHT, \textit{pubsub}, negociação de transporte) torna o sistema difícil de configurar e de entender \cite{irohfaq}.

O Iroh nasceu como um projeto de reescrita do IPFS na linguagem Rust para torná-lo mais rápido. Contudo, para atingir as métricas de desempenho e confiabilidade desejadas, os
desenvolvedores perceberam que precisariam quebrar a compatibilidade com as implementações existentes do IPFS \cite{irohnewdirection}. A motivação do Iroh é fornecer conexões P2P diretas entre dispositivos, identificados por chave pública, com a
menor latência possível e sem a complexidade de configuração da camada de rede \cite{iroh}.

O Iroh foi construído nativamente sobre o protocolo QUIC \cite{rfc9000}, a base do HTTP/3 \cite{rfc9114}, descartando o TCP \cite{iroh}. O QUIC oferece estabelecimento rápido de conexão, criptografia nativa (TLS 1.3) e ausência do bloqueio de cabeça de fila (\textit{head-of-line blocking}) entre \textit{streams},
tornando a troca de mensagens mais rápida, mesmo em redes instáveis, como 4G e 5G.

Diferente de protocolos antigos que usam SHA-256, o iroh-blobs, protocolo de transferência de dados do Iroh, utiliza o algoritmo de hash BLAKE3 \cite{irohblobs}. Por ser internamente uma árvore de Merkle, ele permite que os dados sejam processados de forma paralelizada e verificados durante a transmissão (\textit{verified streaming}), aproveitando a velocidade
máxima do hardware \cite{blake3}.

Ele resolve problemas com NAT (roteadores bloqueando conexões de fora), um dos principais de aplicações P2P, com servidores de repasse (\textit{relays}), originalmente chamados de DERP (\textit{Designated Encrypted Relay for Packets}) \cite{irohchangelog}. A conexão começa passando pelo \textit{relay}, por onde os nós trocam informações de rede e tentam se conectar diretamente (\textit{hole punching}). Se
conseguirem, o \textit{relay} sai do caminho; se os \textit{firewalls} forem muito rígidos, o tráfego continua passando por ele. O \textit{relay} não consegue ler os dados, pois a conexão é criptografada de ponta a ponta \cite{irohrelay}. O usuário não
percebe a diferença, pois ela é abstraída pelo Iroh.

Sobre essa camada de conexão, o Iroh oferece protocolos separados, como o iroh-docs, para sincronização de estados \cite{irohprotocols}. Ele permite que múltiplos nós editem uma estrutura de dados de
chave-valor concorrentemente e a rede cuida de sincronizar as atualizações, através de uma técnica chamada \textit{range-based set reconciliation} \cite{irohdocs, meyer2022}.

O Iroh já é utilizado em aplicações reais. O Delta Chat, um aplicativo de mensagens descentralizado \cite{deltachat}, usa o Iroh em duas funcionalidades: para transferir o
\textit{backup} de um dispositivo para outro através de um QR code \cite{deltachattransfer}, e para criar canais de comunicação em tempo real entre os aplicativos webxdc
que rodam dentro das conversas, através do iroh-gossip \cite{deltachatpeerchannels}. Nos dois casos, os dispositivos se conectam diretamente, sem que os dados passem pelo servidor de mensagens.

Já a Nous Research utiliza o Iroh na Psyche, uma rede para treinamento distribuído de modelos de IA pela internet. Os nós trocam metadados sobre os resultados do treinamento
através do iroh-gossip e baixam os resultados em si através do iroh-blobs. Segundo os autores, o Iroh consegue estabelecer conexões diretas em cerca de 90\% dos casos, contra
70\% da libp2p, e migra a conexão entre redes sem interrompê-la, graças ao QUIC \cite{nouspsyche}.

Por fim, os desenvolvedores do Iroh mantêm uma integração de exemplo com o Automerge, onde o Iroh é usado como camada de rede para sincronizar documentos CRDT entre os nós \cite{irohautomerge}.

\section{Yjs}

A prova acadêmica dos CRDTs resolveram os problemas matemáticos para sincronização de estados, mas consumiam muita memória e CPU no processo. A serialização de grandes JSONs travava navegadores
mais fracos, a depender do tamanho dos documentos serializados. O Yjs foi criado por Kevin Jahns para preencher a lacuna entre a teoria acadêmica e alta performance, implemetando uma biblioteca capaz
de suportar a edição colaborativa em grande escala sem gargalos de memória e CPU.

Yjs implementa o algoritmo YATA, anteriormente desenhado para evitar o embaralhamento de textos quando usuários diferentes editam os mesmos valores. Ele modela o documento colaborativo como uma lista
duplamente encadeada, em vez de árvores. O Run-Length Encoding (RLE) usa codifiação binária compacta para condensar milhares de operações em vetor de poucos bytes. O Yjs é focado nos dados e agnóstico
em relação à rede, permitindo acoplar os deltas binários em WebRTC, WebSockets ou malhas P2P.

O Yjs não faz suposições sobre a rede usada nas implementações. Paradigmas como cliente-servidor são tão usáveis quanto \textit{peer-to-peer}: desde que as mudanças sejam comunicadas, ocorrerá a sincronização dos estados.
Isso ocorre independentemente de a comunicação ser feita por WebSockets, WebRTC (\textit{browsers}) ou outros \textit{providers} \cite{yjs}. Ele suporta edição \textit{offline}, o que permite seguir a filosofia Local-First, onde as alterações são aplicadas
primeiramente no local e depois comunicadas aos outros nodos. Eles, por sua vez, aplicam essas alterações, que são mescladas automaticamente pelos CRDTs, abstraídos como estruturas de dados convencionais do JavaScript.

Internamente, todos os tipos compartilhados (\textit{shared types}) são listas. Isso significa que um único algoritmo, o YATA \cite{nicolaescu2016}, resolve os conflitos de todos os tipos. Cada item tem um ID, formado pelo identificador do cliente e por um relógio de Lamport, e guarda os vizinhos
do momento em que foi inserido. Com isso, dois usuários que inserem no mesmo ponto chegam à mesma ordem final. Um item deletado não desaparece, ele vira um \textit{tombstone} no conjunto de deleções (\textit{delete set}). Para que os metadados não cresçam \textit{ad infinitum},
o Yjs junta em um único item os caracteres digitados em sequência pelo mesmo usuário e descarta o conteúdo dos itens deletados. Em um documento real, com 182 mil inserções e 77 mil deleções, o conjunto de deleções ocupou apenas 4,5 KB \cite{yjsinternals}.

Usado por diversas aplicações \cite{yjs}, o Yjs é utilizado tanto em aplicações cliente-servidor quanto \textit{peer-to-peer}. Entre os usuários mais conhecidos estão o JupyterLab, para notebooks colaborativos, e o JupyterGIS, um editor colaborativo de sistemas de informação geográfica \cite{yjs, jupytergis}. O Room.sh, uma ferramenta de reuniões com funcionalidades embutidas de
desenho, edição e codificação, faz uso do \textit{framework} para sincronizar seus vários estados entre múltiplos usuários. Outras aplicações diferentes também contam com uso do Yjs, como o
Dynaboard, um construtor colaborativo de aplicações web, e o LegendKeeper, uma ferramenta colaborativa para planejamento de campanhas de RPG \cite{yjs}.

\section{Automerge}

A motivação do Automerge era permitir aos programadores a criação de softwares Local-First manipulando estruturas de dados familiares (listas, \textit{strings}, objetos JSON) sem precisar lidar com os problemas complexos de sistemas distribuídos \cite{automerge}. Nascido com foco na experiência do desenvolvedor e na preservação de histórico, o Automerge
funciona como um ``Git para estruturas de dados'' \cite{automerge}. Foi criado pelo laboratório Ink \& Switch, o mesmo que publicou o artigo fundador do Local-First, com participação de Martin Kleppmann, um dos autores desse artigo, e é baseado na pesquisa sobre CRDTs para JSON \cite{kleppmann2019, kleppmannberesford2017}.

Enquanto o Yjs representa internamente todos os tipos como listas, com foco em desempenho, o Automerge modela o documento como uma árvore JSON aninhada \cite{kleppmannberesford2017}. Na versão 3.0, o Automerge passou a usar o seu formato colunar comprimido também durante a execução, e não apenas no armazenamento, reduzindo o uso de memória em mais de 10 vezes \cite{automerge3}.
Ele salva o histórico completo das mutações compactando colunas de dados semelhantes \cite{automerge3}. Já a sincronização envia apenas as alterações que o outro nó ainda não possui \cite{automerge}.

Assim como no Git, cada alteração feita em um documento é registrada como uma \textit{change}, parecida com um \textit{commit}. Cada \textit{change} é identificada pelo
\textit{hash} SHA-256 do seu conteúdo e aponta para as \textit{changes} das quais depende, formando um grafo acíclico dirigido (DAG). As \textit{changes} mais recentes desse grafo,
que ainda não possuem sucessoras, são chamadas de \textit{heads} \cite{automergespec}. Dentro de cada \textit{change}, toda operação recebe um ID formado por um contador e
pelo identificador da réplica (\textit{actor ID}), ou seja, um relógio de Lamport. Isso impõe uma ordem total entre as operações: vence o maior contador e, em caso de empate,
o maior \textit{actor ID} \cite{automergespec}.

Na maioria dos casos, a mesclagem é direta. Alterações em objetos ou propriedades diferentes são simplesmente combinadas. Em listas e textos, todas as inserções e deleções
são preservadas, e inserções concorrentes na mesma posição recebem a mesma ordem em todos os nós \cite{automergeconflicts}. Já para valores numéricos, o Automerge oferece
o tipo \texttt{Counter}, que soma incrementos concorrentes: se dois usuários incrementam um valor 3 ao mesmo tempo, o resultado é 5, e não 4 \cite{automergecounters}.

O único caso sem resolução automática é quando dois usuários alteram a mesma propriedade do mesmo objeto. Nesse caso, o Automerge escolhe um dos valores como vencedor,
de forma determinística e igual em todos os nós, sem depender do relógio físico das máquinas. Os valores perdedores não são descartados: eles ficam acessíveis pela função
\texttt{getConflicts()}, permitindo que a aplicação mostre o conflito ao usuário \cite{automergeconflicts}.

A sincronização entre dois nós é baseada no trabalho de \citeonline{kleppmannhoward2020}. Cada nó guarda o que sabe sobre as \textit{heads} do outro e, a cada mensagem
trocada, envia apenas as \textit{changes} que o outro ainda não possui, até que os dois tenham o mesmo histórico \cite{automergesync}. Por cima disso, o Automerge Repo
cuida da parte de infraestrutura: ele salva os documentos localmente através de adaptadores de armazenamento, como o IndexedDB no navegador ou o sistema de arquivos,
e os distribui para outros nós através de zero ou mais adaptadores de rede, como o WebSocket \cite{automergerepo}.

Os casos de uso podem variar, mas ele se estende para todas as aplicações Local-First ou que requerem um histórico de edições, como no Git. Aplicações Local-First exigem armazenamento local robusto e controle de versão, como gerenciadores de conhecimento, ferramentas de produtividade \textit{offline} e outros tipos de sistemas
onde os usuários necessitam viajar entre os históricos passados de alterações.

Os primeiros casos de uso foram os protótipos do próprio Ink \& Switch, apresentados no artigo sobre Local-First \cite{kleppmann2019}: o Trellis, um quadro Kanban inspirado no
Trello, que experimentou uma interface de ``viagem no tempo'' para navegar pelos estados anteriores do documento; o Pixelpusher, um editor colaborativo de \textit{pixel art}; e o
PushPin, um mural colaborativo de mídias mistas, parecido com o Miro, considerado pelos autores o mais completo dos três. Além disso, como visto na seção sobre o Iroh, existe
uma integração de exemplo que usa o Iroh como camada de rede para sincronizar documentos do Automerge \cite{irohautomerge}.

\section{Comparativo e Posicionamento do Trabalho}

WORK IN PROGRESS - É PRECISO AVANÇAR MAIS PARA CONTINUAR ESSE CAPITULO

\begin{table}[htb]
    \centering
    \caption{Comparativo entre libp2p e Iroh}
    \label{tab:libp2p-iroh}
    \begin{tabular}{|p{0.22\textwidth}|p{0.34\textwidth}|p{0.34\textwidth}|}
        \hline
        \textbf{Critério} & \textbf{libp2p} & \textbf{Iroh} \\ \hline
        Arquitetura & Coleção de módulos empilháveis, conforme a necessidade da aplicação & Núcleo de conexão, com protocolos separados por cima (iroh-blobs, iroh-gossip, iroh-docs) \\ \hline
        Implementações & Várias (Go, Rust, JavaScript, entre outras), seguindo as mesmas especificações & Rust \\ \hline
        Identidade do nó & PeerID, derivado da chave pública & Chave pública Ed25519 \\ \hline
        Transportes & TCP, QUIC, WebRTC, WebSockets, entre outros & QUIC \\ \hline
        Segurança e multiplexação & Módulos próprios sobre a conexão (ou nativos, no QUIC) & Nativos do QUIC \\ \hline
        Travessia de NAT & AutoNAT, Circuit Relay e DCUtR (\textit{hole punching}) & Conexão inicia pelo \textit{relay} e migra para direta via \textit{hole punching} \\ \hline
        Conexões diretas (segundo a Nous Research) & Cerca de 70\% & Cerca de 90\% \\ \hline
        Descoberta de pares & Kademlia DHT, mDNS e \textit{rendezvous} & DNS (Pkarr) e, opcionalmente, a DHT do BitTorrent; sem DHT Kademlia própria \\ \hline
        Difusão de mensagens & GossipSub & iroh-gossip (HyParView e PlumTree) \\ \hline
        Casos de uso & Camada de consenso do Ethereum, IPFS, OrbitDB & Delta Chat, Psyche \\ \hline
    \end{tabular}
    \fonte{Elaborado pelo autor, com base em \citeonline{libp2p}, \citeonline{libp2pspecs}, \citeonline{iroh}, \citeonline{irohfaq}, \citeonline{irohaddresslookup} e \citeonline{nouspsyche}.}
\end{table}

\begin{table}[htb]
    \centering
    \caption{Comparativo entre Yjs e Automerge}
    \label{tab:yjs-automerge}
    \begin{tabular}{|p{0.26\textwidth}|p{0.32\textwidth}|p{0.32\textwidth}|}
        \hline
        \textbf{Critério} & \textbf{Yjs} & \textbf{Automerge} \\ \hline
        Base teórica & YATA \cite{nicolaescu2016} & CRDT para JSON \cite{kleppmannberesford2017} \\ \hline
        Modelo de dados & Todos os tipos são listas internamente & Documento JSON aninhado (mapas, listas, texto, contadores) \\ \hline
        Histórico & Descarta o conteúdo dos itens deletados & Guarda o histórico completo, como um grafo de \textit{changes} \\ \hline
        Foco & Edição colaborativa em tempo real e desempenho & Versionamento e aplicações Local-First \\ \hline
        Rede & Agnóstico, via \textit{providers} & Agnóstico, via Automerge Repo \\ \hline
        Tempo para aplicar o histórico real (B4) & 5.714 ms & 14.326 ms \\ \hline
        Tamanho final do documento (B4) & 159.929 bytes & 129.116 bytes \\ \hline
    \end{tabular}
    \fonte{Elaborado pelo autor, com base em \citeonline{yjs}, \citeonline{yjsinternals}, \citeonline{automerge} e \citeonline{crdtbenchmarks}.}
\end{table}
